\section{Introduction}
\label{sec:intro}

Enterprise AI systems increasingly rely on multi-agent architectures, where complex tasks are decomposed across chains of specialized language models. A growing empirical literature documents the performance of these systems---yet the fundamental trade-off governing their optimal depth remains theoretically unexplored. As context windows expand by orders of magnitude (from GPT-3's 2K tokens to Gemini-1.5's 1M tokens) and API costs continue to decline, the natural intuition is that deeper chains---more layers of specialized processing---should become optimal. We show that this intuition is incomplete: even with costless communication, endogenous constraints on information transmission bound the efficient depth of the chain.

The core economic mechanism is \emph{information depreciation}---the gradual loss of signal quality as information traverses a processing hierarchy. In a multi-agent chain, each layer receives noisy signals from the layer below, updates beliefs, and transmits a summary upward. Because each agent's context window is finite, not all incoming signals can be attended to with full precision; some are compressed, some are discarded, and the resulting transmission is necessarily coarser than the input. This degradation compounds across layers, so that the signal reaching the terminal layer carries only a fraction of the original information content. Information depreciation thus operates analogously to capital depreciation: each processing layer is an investment in refined decision-making, but that investment suffers a depreciative loss before it reaches the next stage.

We formalize this intuition by embedding rational inattention \citep{sims2003implications} into a multi-layer information-processing hierarchy. The model delivers two core structural results. First, the optimal budget allocation across layers is \emph{front-loaded}: the largest attention budget should be deployed at the earliest layer, with monotonically decreasing allocations at higher layers (Proposition~\ref{prop:front_loading}). This is the model's most empirically testable prediction---it can be evaluated using only observable architectural choices and does not require calibrating the unobserved depreciation parameter. In practice, most production multi-agent systems already use this architecture: the largest model handles the initial query understanding and retrieval, while smaller, faster models handle intermediate reasoning steps. Our model provides the first theoretical justification for this engineering practice.

Second, two conditions jointly generate finite optimal depth. Signals depreciate in quality as they traverse the chain, formalized as an \emph{endogenous} retention factor $\eta_\ell \in (\underline{\eta}, \bar{\eta})$ that captures the natural decay of information during processing at layer $\ell$ (Proposition~\ref{prop:depth_finite}). The preservation rate $\eta_\ell$ is determined by the average attention allocated per signal at layer $\ell$ through Assumption~\ref{ass:endogenous_eta}, with the key property that even unlimited attention cannot push preservation above a technological bound $\bar{\eta} < 1$. Each agent also faces a finite Shannon-capacity budget for attention allocation, which determines \emph{where} in the interval $[\underline{\eta}, \bar{\eta}]$ the actual preservation rate lies. The depth of the chain is bounded because the compounding of these two forces---inherent algorithmic noise (captured by $\bar{\eta} < 1$) and endogenous attention loss---eventually renders marginal layers unproductive. Either condition alone is insufficient: with $\bar{\eta}=1$ and no inherent degradation, information could circulate without bound regardless of attention budgets; without the bound $\bar{\eta} < 1$, unlimited attention would suffice to maintain signal quality at any depth.

Our analysis is related to the canonical models of multi-layer organization developed by \citet{garicano2000hierarchies} and \citet{dessein2006adaptive} (henceforth, GHM). In the GHM framework, finite organizational depth arises from the interplay of diminishing returns and increasing marginal costs of communication or effort. While information loss in hierarchies is a familiar theme in organizational economics \citep[e.g.,][]{sah1991fallibility,dessein2002authority,qian1994incentives}, existing models typically treat communication fidelity as an exogenous parameter or as the outcome of explicit incentive problems. Our model \emph{complements} the GHM framework by adding a micro-foundation for information processing: the transmission factor $\rho_\ell$, defined as the ratio of achieved posterior precision to first-best posterior precision, emerges endogenously from the optimal attention-allocation problem and reflects the actual information degradation at each layer. This connects the organizational model to the Bayesian micro-foundations of rational inattention \citep{sims2003implications} and generates predictions about how attention weights, signal heterogeneity, and capacity constraints jointly determine organizational form.

This paper makes three contributions. First, we introduce the concept of \emph{information depreciation} into the theory of organizational hierarchies. By modeling signal quality as a depreciating asset that traverses a processing chain, we provide a tractable framework that links communication-loss models in organizational economics \citep{sah1991fallibility,dessein2002authority} to the rational-inattention literature \citep{sims2003implications}. We \emph{endogenize} the depreciation rate by making it depend on the attention allocation through Assumption~\ref{ass:endogenous_eta}, so that the preservation parameters $(\underline{\eta}, \bar{\eta}, a)$ discipline the model's predictions and can be calibrated from observable features of LLM systems.

Second, we provide a mapping from LLM architecture parameters---context window length, token-level precision, and API pricing---to predictions about optimal organizational depth and budget allocation. The model generates sharp comparative statics: depth increases with context window size and decreases with signal depreciation; the optimal allocation is front-loaded regardless of cost structure; and API costs affect the intensive margin of processing but are largely irrelevant for the extensive margin of depth once costs are below a threshold. The front-loading prediction is especially amenable to empirical testing because it requires only observable architectural data---the model sizes or context-window lengths assigned to each layer---and can guide system-design decisions in enterprise AI deployment.

Third, we calibrate the model using publicly available data on benchmark accuracy, context window capacities, and API pricing for state-of-the-art language models. The calibration demonstrates that the key predictions hold for empirically plausible parameter values and generates quantitative predictions about the range of efficient depth under current and projected technology. We provide comparative statics that can discipline future empirical work on multi-agent AI architectures.

The remainder of the paper is organized as follows. Section~\ref{sec:literature} reviews the related literature. Section~\ref{sec:model} presents the model. Section~\ref{sec:single} analyzes the single-layer attention-allocation problem and derives the endogenous transmission factor $\rho_\ell$. Section~\ref{sec:multi} extends the analysis to multiple layers, establishes that the optimal budget allocation is front-loaded, and characterizes the finite optimal depth. Section~\ref{sec:calibration} calibrates the model using LLM technology data and presents comparative statics. Section~\ref{sec:empirics} discusses testable implications for the design of multi-agent systems. Section~\ref{sec:extensions} considers extensions, and Section~\ref{sec:conclusion} concludes.
